\section{The stabilized obstruction from the torical theorem}
\label{sec:obstruction}

The shortening construction varies the number of circle factors while
keeping the base dimension fixed. We prove the required obstruction with
a threshold independent of that circle count. Its only external geometric
input is the local torical theorem of Cecchini and Schick; the reduction
below checks its compactness, degree and local metric hypotheses.

\begin{definition}
\label{def:torical-metric}
A \emph{torus-stabilized metric} on a smooth manifold $Y$ is a metric
\begin{equation}
\label{eq:stabilized-metric}
 \mathcal G=g+\sum_{\nu=1}^{b}m_\nu^2\,d\theta_\nu^2
 \qquad\text{on }Y\times\T^b,
\end{equation}
where $b\geq0$ is finite, $g$ is a smooth metric on $Y$, and the functions
$m_\nu:Y\to(0,\infty)$ are smooth. Each circle is
$\R/(2\pi\Z)$. We call $g$ the base metric.
Scalar curvature always refers to the full metric $\mathcal G$.
Differentials of base maps have the norm determined by $g$.
\end{definition}

For a map $F:Y^d\to S^d$ equal to a fixed value $y_\infty$ outside a
compact set, degree means evaluation against a smooth top-degree form
supported away from $y_\infty$ and of integral one. Equivalently, it is the
signed preimage count at a regular value different from $y_\infty$.
This convention also applies when $Y$ is disconnected.

\begin{theorem}[Stabilized torical obstruction]
\label{thm:torical-obstruction}
For every integer $d\geq2$ and every $\chi>0$ there is a number
$\delta(d,\chi)>0$ for which the following data do not exist:
\begin{enumerate}
 \item a smooth oriented $d$-manifold $Y$ without boundary, possibly
 disconnected, with a torus-stabilized metric $\mathcal G$ satisfying
 $\Scal_{\mathcal G}\geq\chi$;
 \item a smooth map $F:Y\to S^d$, equal to one value outside a compact
 set, such that $\deg F\ne0$ and
 $\norm{dF}_{\mathrm{op}}\leq\delta(d,\chi)$ everywhere.
\end{enumerate}
The constant is independent of the finite circle count in
\cref{eq:stabilized-metric}. Neither the base metric nor the stabilized
metric is assumed complete.
\end{theorem}

Cecchini--Schick prove that if a closed oriented Riemannian $D$-manifold $X$
has a smooth map $f:X\to S^1\times\T^{D-1}$ of nonzero degree, and a
connected open arc $J\subset S^1$ satisfies
\[
 \Scal_X\ge\chi>0,\qquad |df|_{\mathrm{op}}\le\epsilon
 \quad\text{on }f^{-1}(J\times\T^{D-1}),
\]
then
\begin{equation}\label{eq:cecchini-schick}
 2\pi\epsilon\ge\sqrt\chi\,\operatorname{Length}(J).
\end{equation}
This is \cite[Theorem~1.14(1)]{CecchiniSchick2021}. The statement has no
upper dimension bound and no spin hypothesis. Its notion of contraction is
a pointwise bound on the full differential. In high dimensions its proof uses the
minimal-slicing input identified in that paper; we invoke the
published theorem at this interface. We use its flat torus normalization
in the target; any fixed change from our angular convention is absorbed
in the tube-projection bound and the constant circle enlargements below.


\begin{proof}[Proof of \cref{thm:torical-obstruction}]
Fix $d\ge2$ and $\chi>0$. Suppose that the data in the theorem exist,
and write $b$ for the finite number of circles. Set $D=d+b$.
Thus $Y^d$ carries the stabilized metric
\[
 g+\sum_{\nu=1}^b m_\nu^2d\theta_\nu^2
\]
of scalar curvature at least $\chi$, and the nonzero-degree map
$F:Y\to S^d$ has differential bounded by a number $\delta>0$ in the
base metric $g$. We derive a positive lower bound for $\delta$ which
is independent of $b$.

Choose a fixed two-sided embedded $\T^{d-1}$ in $S^d$ away from the exterior
value of $F$, with a closed tubular neighborhood parametrized by
$\T^{d-1}\times[-2,2]$. Such embeddings can be built inductively in a
Euclidean chart: add one ambient dimension and take the boundary of a small
trivial two-dimensional normal disk bundle around the preceding torus.
After a rotation, the target tube and all its projection bounds depend only
on $d$.

Choose regular interval levels
$s_-\in(-5/4,-3/4)$ and $s_+\in(3/4,5/4)$ and pull back the strip
between them. This gives a compact smooth band $W\subset Y$, with a
map to that strip of nonzero relative degree. Compactness follows because
the closed target tube avoids the exterior value. Regularity of the two
levels gives a smooth boundary; nonzero relative degree follows by
counting preimages of a regular interior target of $F$.

Write the target strip as $\T^{d-1}\times[s_-,s_+]$. Double
$W$ with the opposite orientation on the second copy. On both copies
retain the $\T^{d-1}$ component of the target map, and send the interval
component monotonically along complementary semicircles of $S^1$.
To make this map smooth across the seams, first choose disjoint source
boundary collars and precompose the \emph{entire base map}, including its
torus component, with a smooth collar reparametrization that is constant
in the normal variable near the boundary. On the first copy, choose the circle angle $\alpha(t)$ to be
$\pi/2+\pi t/2$ for $|t|\le1/2$, extend it monotonically to $0$ and
$\pi$ at the two walls, and make it flat near the endpoints. On the
second copy use the complementary semicircle, with the same endpoint
values in $S^1$. The two torus components then
agree and are constant in the normal variable at each seam, while the
circle components agree there as well. Hence they define a smooth map
\[
 f:\operatorname{dbl}(W)\times\T^b
       \longrightarrow S^1\times\T^{d-1}\times\T^b,
\]
with the identity on the final torus.

All collar changes can be confined near the walls, away from $|t|\le1/2$.
Choose the open target arc
\[
 J=\{e^{i\alpha}:3\pi/8<\alpha<5\pi/8\}.
\]
Its full preimage in the first copy is precisely $|t|<1/4$, where the
map is unchanged and its interval slope is the fixed number $\pi/2$.
The second copy maps to the complementary semicircle and has no preimage
over $J$. Counting preimages of a regular target over $J$ therefore
recovers the nonzero relative degree of the original band map. The map
$f$ has nonzero degree, on at least one connected component if the
double is disconnected.

Each $m_\nu$ has a positive minimum on the compact band. Enlarge its circle
length by a constant factor until the joint projection to the standard
$\T^b$ has operator norm at most $\delta$ on $W$. This changes no local
scalar curvature: a constant change of a circle coordinate is a local
isometry between the two warped metrics. The operator norm of the joint
projection is the maximum inverse circle length, rather than a sum over
$b$.

Give the closed double any smooth metric agreeing with the enlarged
stabilized metric throughout the chosen middle substrip and a neighborhood
of its closure. On the preimage of $J\times\T^{d+b-1}$ it then has scalar
curvature at least $\chi$. The base part of $f$ has differential at most
$C_d\delta$, because only the fixed tube projection and a fixed middle-arc
parametrization are used there. Its torus part has differential at most
$\delta$. The orthogonal block structure therefore gives
\[
 |df|_{\mathrm{op}}\le\max(C_d,1)\delta
 \quad\text{on }f^{-1}(J\times\T^{d+b-1}),
\]
independently of $b$. No metric or curvature control elsewhere on the
double is required by the cited theorem.

Applying \eqref{eq:cecchini-schick} gives
\[
 2\pi\max(C_d,1)\delta
 \ge\sqrt\chi\,\operatorname{Length}(J).
\]
The fixed tube and middle-arc choices depend only on $d$. In particular
we may take
\[
 \delta(d,\chi)=
 \frac{\sqrt\chi\,\operatorname{Length}(J)}
      {4\pi\max(C_d,1)}.
\]
A map with differential at most this number contradicts the preceding
lower bound. The threshold depends only on $d$ and $\chi$, independently
of the finite circle count $b$.
\end{proof}
